%==================================================
% ----- Mapa de estado e saídas cinemáticas
%==================================================
\section{Mapa de estado e saídas cinemáticas}

O estado integrado é
$\vec x=\begin{bmatrix}\vec q\\ \dot{\vec q}\end{bmatrix}$,
em que $\vec q=\begin{bmatrix} q_b \\ \vec q_m \end{bmatrix}
$ empilha a atitude da base (quaternion $q_b$) e as coordenadas das juntas do manipulador, e
$\dot{\vec q}=\begin{bmatrix} \boldsymbol\omega_b \\ \dot{\vec q}_m \end{bmatrix}$
empilha as respectivas velocidades. 
A consistência do quaternion ($q_b^{\top}q_b=1$)
é a restrição holonômica apresentada anteriormente
\fn{Ver \eqref{eq:PhiJ1} e \eqref{eq:PhiJ2}}.

Para fins de apresentação e comparação, adotou-se: 
\begin{equation}
Y_1 \equiv \vec q,\qquad
Y_3 \equiv \dot{\vec q},\qquad
Y_2 \equiv h(\vec q),\qquad
Y_4 \equiv J_h(\vec q)\,\dot{\vec q},
\label{eq:mapa_estado}
\end{equation}
sendo que $h(\vec q)$ reúne as grandezas geométricas de interesse e uma parametrização mínima da orientação do efetuador ou o próprio
${}^{B}\!q_E(\vec q)$,
e $J_h(\vec q)=\partial h/\partial \vec q$
é a jacobiana correspondente. 
Assim, $Y_2$ e $Y_4$ são saídas algébricas calculadas a partir de $(Y_1,Y_3)$ via cinemática direta e jacobianas,
conservando-se no integrador apenas $\vec x$. 
Os vínculos e seus níveis diferencial e de aceleração permanecem impostos por
\eqref{eq:vellevel} e \eqref{eq:acclevel},
bem como pela estabilização de Baumgarte quando empregada
\eqref{eq:acclevel_baumgarte}.


